\section{The Expected Discovery Metric}\label{sec:framework}

% Define the metric
To quantify gene discovery as a continuous function of statistical power, we define an \textit{expected discovery} metric that measures how much evidence has accumulated for each gene above a baseline prior. For any evidence type (direct, indirect, or combined), the expected discovery for gene $i$ is the posterior probability of disease relevance minus the prior probability: $D_i = P_i^{\text{post}} - P_i^{\text{prior}}$, where $P_i^{\text{post}} = \pi \cdot \text{ABF}_i / (1 + \pi \cdot \text{ABF}_i)$ and $\pi = 0.05$ is the background prior (Methods). The total expected discovery $\tilde{N} = \sum_i D_i$ represents the expected number of genes identified above baseline---a continuous metric that avoids arbitrary significance thresholds while remaining interpretable as a gene count.

% Downsampling procedure
To trace how gene discovery accumulates with statistical power, we applied a downsampling procedure to {{NUM_TRAITS_SATURATION_ANALYSIS}} common traits with meta-analyzed GWAS summary statistics. For each trait, we simulated reduced sample sizes by inflating the standard error of each SNP's effect estimate according to a downsampled fraction $f \in \{0.1, 0.2, \ldots, 1.0\}$ (Methods). At each fraction, we computed expected discovery separately for direct support ($\tilde{N}_D$), indirect support ($\tilde{N}_I$), and combined support ($\tilde{N}_C$), yielding discovery trajectories as a function of effective sample size.

% Classification approach
We classified each trait's discovery status by comparing the fit of two models to its direct support trajectory. A linear model implies that gene discovery continues to increase proportionally with sample size---the trait remains power-limited. An exponential saturation model implies that discovery is approaching an asymptote---the trait has reached or is entering saturation. We compared models using the Bayesian Information Criterion and further quantified the degree of saturation as the proximity to the fitted asymptote (Methods). Results were robust to the choice of background prior (Supplementary Figure~\ref{fig:prior-sensitivity}).
